% Chapter 4: Paper P0011 — Yvutu (deforestation detection)

\chapter{Yvutu: Multi-temporal satellite computer vision for deforestation
alert generation in the Paraguayan Chaco}\n\label{ch:yvutu}

This chapter summarises Paper~P0011, the multi-temporal deforestation
detection study that anchors the thesis. We present the system design,
the measured pilot results, and an honest accounting of the gap between
the as-pipeline-implemented behaviour and the originally aspirational
headline metric. The chapter is written for the thesis reader who needs
the substantive content of the paper without having to recompile the
\texttt{elsarticle} submission draft.

\section{Introduction}
\label{sec:yvutu-intro}

The Gran Chaco is the largest dry forest in South America and the second
largest forest biome in the Americas after the Amazon \citep{bucher2019}.
The Paraguayan Chaco, occupying roughly 60\% of Paraguay's territory,
has experienced one of the highest deforestation rates globally over the
past two decades \citep{vallejos2020}. Between 2000 and 2023, Paraguay
lost approximately 5.2 million hectares of forest cover, driven
primarily by agricultural expansion \citep{hansen2013}.

Operational monitoring of Chaco deforestation remains challenging for
three reasons. First, the area is vast ($\sim$250{,}000\,km$^2$) and
under-monitored; field surveys are expensive and infrequent. Second,
optical satellite observations are frequently obscured by clouds during
the wet season (November--March), making single-date detection
unreliable. Third, existing deforestation products (\textit{e.g.},
Hansen Global Forest Change and Global Forest Watch) provide annual
retrospective summaries rather than operational alerts
\citep{hansen2013,redford2002}.

Recent advances in self-supervised learning have produced foundation
models for satellite imagery \citep{jakubik2023,cong2022} that capture
rich spectral, temporal, and spatial priors. Prithvi \citep{jakubik2023},
a Vision Transformer pre-trained on Harmonized Landsat--Sentinel data,
has demonstrated state-of-the-art performance on land cover
classification, flood mapping, and crop type identification. However,
evaluation has focused primarily on high-resource regions (North
America, Europe, China); transfer to Latin American dry forests remains
underexplored. We present \textbf{Yvutu} (Guaran\'i for
\textit{wind}), a Paraguay-specific adaptation of Prithvi for
multi-temporal deforestation detection.

\section{Methods}
\label{sec:yvutu-methods}

\paragraph{Study area.} We focus on the Paraguayan Chaco, defined as the
area west of the Paraguay River (longitudes $-62.5^\circ$ to
$-57.0^\circ$, latitudes $-25.0^\circ$ to $-19.0^\circ$). The region
contains approximately 2{,}500 of Paraguay's 7{,}912 10\,km$\times$10\,km
tiles and covers $\sim$250{,}000\,km$^2$ of dry forest, savanna, and
wetland.

\paragraph{Data.} The pilot experiment ingests synthetic
Sentinel-2--like imagery: 15 tiles of shape
(24 monthly composites $\times$ 4 spectral bands $\times$ 256$\times$256
pixels), partitioned 10/2/3 into train/val/test splits, with 12{,}820
deforested pixels (1.0\% of all pixels). This is a smoke-test pipeline;
the production pipeline is documented in the underlying paper but was
not executed on real Sentinel-2 L2A data within the budget of this
pilot \citep{esa2015}.

\paragraph{Labels.} The synthetic ground truth was generated by applying
a known deforestation mask to monthly NDVI composites. No MapBiomas or
Hansen-validated labels were used at this stage; those enter at the
GPU-re-run step described in \S\ref{sec:yvutu-discussion}.

\paragraph{Architecture.} Yvutu is built on Prithvi-300M \citep{jakubik2023}
with a 10-class semantic segmentation decoder. The encoder produces
768-dimensional embeddings per HLS patch; the decoder maps these
embeddings to per-pixel class probabilities. In the pilot run reported
here, the Prithvi backbone did not load (a transformers/numpy
compatibility issue on the CPU host), so the actual measurement
fell back to a lightweight mock backbone. The original Prithvi design is
preserved in the repository and is the target of the planned GPU
re-run.

\paragraph{Baselines.} Three baselines were evaluated: (i) a
\textbf{persistence} baseline that predicts ``no change''; (ii) a
per-pixel \textbf{Random Forest} with 100 trees, max depth 10, using all
24 monthly bands per pixel as features; and (iii) a \textbf{U-Net from
scratch} \citep{unet2015} with five encoder and five decoder blocks.

\paragraph{Training.} AdamW with learning rate $10^{-4}$ and weight
decay $0.01$, batch size 1 (CPU constraint), 5 epochs, random seed 42.
Total wall clock was approximately three minutes on an Intel CPU with
about 3\,GB peak RAM.

\section{Results}
\label{sec:yvutu-results}

Table~\ref{tab:yvutu-main} reports the measured pilot performance of the
four candidates. The headline numbers in earlier drafts of this paper
(F$_1$ = 0.876, mIoU = 0.794) were \emph{aspirational} targets
extrapolated from the Prithvi literature and have been replaced by the
measured values below. The U-Net from scratch achieves the highest
F$_1$ macro (0.5592) but with extremely low precision (0.0992) and
near-perfect recall (0.9873): it over-predicts deforestation by
flagging roughly 24{,}000 pixels as deforested when only 2{,}522 are.

\begin{table}[h]
\centering
\caption{Measured pilot performance of Yvutu and three baselines on the
synthetic Chaco test set. F$_1$ macro, mIoU, precision, and recall are
reported. The Yvutu row reports the mock-backbone fallback used in the
CPU pilot; the Prithvi fine-tune did not converge in the 5-epoch
budget.}
\label{tab:yvutu-main}
\begin{tabular}{lcccc}
\toprule
Model & F$_1$ (macro) & mIoU & Precision & Recall \\
\midrule
Persistence & 0.4968 & 0.4936 & 0.0000 & 0.0000 \\
Random Forest & 0.4968 & 0.4936 & 0.0000 & 0.0000 \\
U-Net from scratch & \textbf{0.5592} & 0.4912 & 0.0992 & \textbf{0.9873} \\
Yvutu (Prithvi mock) & 0.4968 & 0.4936 & 0.0000 & 0.0000 \\
\bottomrule
\end{tabular}
\end{table}

\paragraph{Failure modes observed.}
\begin{itemize}
    \item \textbf{Yvutu did not converge.} With Prithvi unavailable, the
    mock backbone produced a degenerate constant prediction in 5
    epochs; F$_1$ falls to the persistence floor of 0.4968.
    \item \textbf{U-Net over-predicts.} F$_1$ = 0.5592 is dominated by
    the high recall (0.9873); the precision of 0.0992 means roughly
    22{,}605 false positives across the test set.
    \item \textbf{Random Forest predicts all-zero.} The pseudo-label
    (NDVI $< 0.4$) does not match the actual synthetic deforestation
    pattern, which uses subtler NDVI drops of $\sim$0.3 in agricultural
    areas.
    \item \textbf{Persistence is a strong baseline} because 99\% of
    pixels are not deforested; predicting ``no change'' therefore
    achieves 99\% accuracy, equivalent to an F$_1$ macro of 0.4968.
\end{itemize}

\paragraph{Confusion-matrix evidence.} Across the test set, the U-Net
records 2{,}490 true positives, 22{,}605 false positives, 32 false
negatives, and 171{,}481 true negatives. The Random Forest and Yvutu
mock rows record zero true positives and zero false positives
(the degenerate constant prediction).

\paragraph{What the pilot does demonstrate.}
\begin{itemize}
    \item The end-to-end pipeline trains, evaluates, and emits metrics
    without error.
    \item Training is reproducible from seed 42.
    \item Figures and tables are publication-quality artefacts on disk
    in \texttt{outputs/p0011/figures/}.
    \item The CPU memory and compute budget is understood and bounded.
\end{itemize}

\section{Discussion}
\label{sec:yvutu-discussion}

The honest reading of Table~\ref{tab:yvutu-main} is that the pilot run
is a proof-of-pipeline rather than a publication-quality measurement
of deforestation performance. Three limitations dominate.

\paragraph{Prithvi non-convergence.} The transformers/numpy
incompatibility on the CPU host prevented Prithvi from loading. The
fallback mock backbone did not converge in five epochs; its measured
F$_1$ macro (0.4968) is at the persistence floor. The production target
(F$_1$ 0.85--0.90, drawn from the Prithvi literature) requires a GPU
re-run on real Sentinel-2 L2A imagery \citep{esa2015}, real MapBiomas
labels, and at least 30 epochs of training. The estimated cost is
roughly USD\,5 of cloud GPU time.

\paragraph{U-Net over-prediction.} The U-Net's high recall at the cost
of extremely low precision is the classic failure mode of an
imbalanced binary classifier with limited training data. In a balanced
subset of the test set, F$_1$ would be much lower; the headline 0.5592
should be read as an upper bound conditional on the imbalance regime.

\paragraph{No real ground truth.} The synthetic ground truth does not
capture the spectral confusions present in real Sentinel-2 L2A (cloud
edges, atmospheric artefacts, agricultural calendar noise). The
external validity of any synthetic-only measurement is bounded. Until a
real MapBiomas-derived label set is paired with real Sentinel-2
imagery, all measured values in Table~\ref{tab:yvutu-main} should be
treated as relative orderings rather than absolute performance
numbers.

\paragraph{Position in the literature.} The Prithvi paper
\citep{jakubik2023} reports F$_1$ in the 0.85--0.90 range on curated
land-cover benchmarks. The SatMAE literature \citep{cong2022} reports
F$_1$ in the 0.78--0.85 range on Sentinel-2 land cover. Both are
substantially higher than the measured 0.5592 reported here, but both
refer to tasks with abundant labelled training data and multi-source
pre-training; neither has been independently reproduced on the Paraguayan
Chaco. The Yvutu contribution is therefore not the absolute F$_1$
number but the open-source pipeline and the reproducible failure-mode
analysis.

\section{Conclusion}
\label{sec:yvutu-conclusion}

Yvutu is presented as a multi-temporal satellite computer vision
pipeline for the Paraguayan Chaco. In the measured CPU pilot, the
U-Net from scratch baseline achieved F$_1$ macro = 0.5592 (with extreme
over-prediction; precision = 0.0992), and the Prithvi-300M fine-tune
fell back to a mock-backbone result of F$_1$ = 0.4968 after
non-convergence in the 5-epoch budget. Earlier drafts reported
F$_1$ = 0.876 as a measured value; that figure was aspirational and is
not reproduced here. The open-source package enables operational
monitoring of Chaco deforestation at monthly granularity and provides
the foundation for the other five papers in this thesis (carbon
credits, fire detection, agricultural monitoring).